\begin{abstract}
This paper describes a Python environment for deferred-time granular synthesis. 
The system does not synthesise grains from a Gabor primitive but granulates recorded 
material: it belongs to what Lippe terms granular sampling, where the dominant 
expressive parameter is the read position in the source buffer. Granulating at musical 
densities requires specifying thousands of grains per second; the system's problem 
is explicit, legible parametric control over that mass of data. Parameters are specified 
in a declarative approach with a YAML domain-specific language — validated during writing by a 
dedicated Language Server — and interpreted as tendency masks: a central trajectory, 
a deviation range, and an independent per-grain sampling. A decoupled renderer 
architecture supports a per-stream STEMS workflow whose incremental cache and 
DAW-project export keep the edit–listen cycle short enough to inhabit. An 
automatically generated graphical score, whose vertical axis maps the read 
position rather than frequency, closes the feedback loop as a read-only 
diagnostic. We situate the system within the CIM granular tradition and 
discuss what is at stake in choosing deferred time while real time is available.
\end{abstract}
